\chapter{Simulation}
\label{ch:sim}

\section{Launching}

\begin{lstlisting}
# headless server, the default for scripted checks
ros2 launch baxter_gz_sim sim.launch.py headless:=true

# with the Gazebo GUI
ros2 launch baxter_gz_sim sim.launch.py headless:=false

# with RViz as well
ros2 launch baxter_gz_sim sim_rviz.launch.py headless:=false
\end{lstlisting}

Run one simulation at a time.

\begin{notebox}[Do not reach for \texttt{ROS\_DOMAIN\_ID}]
It is tempting to isolate concurrent runs with a domain ID or a Gazebo partition.
Do not. The documented workflow runs a single simulation, and setting a domain has
silently split the bridge from the action shim --- two halves of one system, each
convinced the other was absent. If you need isolation, stop the other simulation.
\end{notebox}

\section{The model}

The simulation URDF fixes \texttt{world}~$\rightarrow$~\texttt{base} at
$z=0.92418$\,m, which puts the bottom of the pedestal on the floor. This is a
fixed joint, not a floating base with a settle phase: the robot is bolted down in
reality and a drop-and-settle would only add nondeterminism to every test that
follows. Both arms start at the SRDF neutral state.

\begin{figure}[h]
\centering
\graphicsorplaceholder{figures/sim_rviz.png}{\linewidth}{the plain RViz profile
with the RobotModel and TF displays}
\caption{\texttt{sim\_rviz.launch.py}. Fixed frame \texttt{world}, Global Status
Ok, RobotModel and TF displays active. This is the profile to check first: if the
model does not appear here, nothing in Chapter~\ref{ch:moveit} will work either.}
\label{fig:simrviz}
\end{figure}

\begin{figure}[h]
\centering
\graphicsorplaceholder{figures/tf_tree.pdf}{\linewidth}{the kinematic tree
derived from the simulation URDF}
\caption{Kinematic tree from \texttt{baxter\_gz\_control.urdf.xacro}. Dashed
edges are fixed joints; the arm joints are the seven per side that carry position
commands.}
\label{fig:tftree}
\end{figure}

Seventeen joints publish state: fourteen arm joints, \texttt{head\_pan}, and one
finger joint per gripper. Fourteen accept position commands --- the arms only.

\section{Controllers}

\begin{table}[h]
\centering
\small
\begin{tabularx}{\linewidth}{@{}lX@{}}
\toprule
\textbf{Controller} & \textbf{Interface} \\
\midrule
\texttt{left\_arm\_controller} & \texttt{/left\_arm\_controller/follow\_joint\_trajectory} \\
\texttt{right\_arm\_controller} & \texttt{/right\_arm\_controller/follow\_joint\_trajectory} \\
\texttt{joint\_state\_broadcaster} & publishes \texttt{/joint\_states} \\
\bottomrule
\end{tabularx}
\end{table}

\begin{lstlisting}
ros2 control list_controllers      # all three must be active
ros2 topic echo /joint_states --once
\end{lstlisting}

Command limits are enforced at the controller manager
(\texttt{enforce\_command\_limits: true}), and every arm joint carries a
trajectory and goal constraint. The goal constraint is $0.02$\,rad, which is the
number every example client checks itself against.

\section{The smoke client}

\begin{lstlisting}
ros2 launch baxter_examples sim_tiny_trajectory.launch.py
\end{lstlisting}

For each arm in turn the client reads a fresh joint state, picks a limit-safe
reversible target, commands it, verifies the final position against a fresh
reading, and returns to the measured start. Action success alone is not treated
as a pass: the client re-reads the state and compares. An action server that
reports success while the arm did not move is exactly the failure worth catching.

Useful parameters:

\begin{table}[h]
\centering
\small
\begin{tabularx}{\linewidth}{@{}llX@{}}
\toprule
\textbf{Parameter} & \textbf{Default} & \textbf{Effect} \\
\midrule
\texttt{joint} & \texttt{s1} & Which joint moves. Any of \texttt{s0 s1 e0 e1 w0 w1 w2}. \\
\texttt{offset} & 0.35\,rad & How far. \\
\texttt{duration} & 3.0\,s & How long, so offset/duration sets speed. \\
\texttt{cancel\_after\_sec} & 0.0 & Cancel mid-motion and verify the hold. \\
\bottomrule
\end{tabularx}
\end{table}

Both arms always move the same joint. That is the point of the parameter: on
hardware, comparing the same joint left against right is what distinguishes an
asymmetric \emph{plan} from an asymmetric \emph{arm}
(Section~\ref{sec:moveit-hardware}).

\section{Cancellation}

\begin{lstlisting}
ros2 run baxter_examples sim_tiny_trajectory --ros-args \
    -p use_sim_time:=true -p cancel_after_sec:=1.0
\end{lstlisting}

The requirement is not merely that cancellation is accepted. The arm must
\emph{hold} where it stopped, not sag, and the client verifies that by measuring
drift after the cancel. A cancel that drops the arm into gravity compensation
would pass a naive check and is a safety-relevant difference on hardware.

\section{Expected diagnostics}

Some Gazebo complaints are normal and should not be chased:

\begin{table}[h]
\centering
\small
\begin{tabularx}{\linewidth}{@{}lX@{}}
\toprule
\textbf{Message} & \textbf{Why it is expected} \\
\midrule
DART cannot create the gripper mimic constraint & Arm-only scope. Gripper TF is available; contact fidelity is not claimed. \\
Controller period 100\,Hz slower than physics at 1\,kHz & Controllers need not run every physics step. \\
Executor unavailable during hardware initialization & Benign if successful initialisation follows immediately. \\
Gazebo exits $-2$ after one Ctrl+C & Normal SIGINT status. \\
\bottomrule
\end{tabularx}
\end{table}

Missing joints, disabled command limits, controller failure, TF gaps, hangs, and
leftover processes are \emph{not} in this category.

\section{Scope}

Only arm joints are controlled. Grippers, cameras, obstacle sensing, and
compatibility fallbacks are outside this profile. Headless Gazebo is camera-less
by construction, so any camera work needs a dedicated path that does not exist
yet.
